(1) 自旋半体系瞬时哈密顿量 \(H_0(t)=-(\hbar\Omega/2)\boldsymbol n(t)\cdot\boldsymbol\sigma\)，\(\Omega>0\) 恒定，单位矢量 \(\boldsymbol n=(\sin\theta\cos\varphi,\sin\theta\sin\varphi,\cos\theta)\)。不作慢变近似，利用Bloch矢量运动方程构造辅助控制哈密顿量 \(H_{\rm cd}(t)\)，使初始处于 \(H_0\) 基态的自旋精确跟随 \(\boldsymbol n(t)\)。
(2) 当 \(\varphi=0\) 而 \(\theta(t)\) 任意光滑变化时，写出显式 \(H_{\rm cd}\) 及全时间传播子 \(U(t,0)\)，验证其解决薛定谔方程并给出对任意瞬时两本征态均无跃迁的理由。
(3) 对一般 \(\theta(t),\varphi(t)\)，计算控制场的瞬时算符范数、被跟踪基态上的期望及方差；说明该场令态跟随不等于控制能量尺度为零，给出 \(\dot{\boldsymbol n}=0\) 的极限。
(4) 若 \(\theta=\theta_0\in(0,\pi)\)、\(\varphi(t)=\omega t\)，经过一个方向闭路 \(T=2\pi/|\omega|\)，求精确跟随基态所得的动力学相位和几何相位（规定 \(\omega>0\)），并说明该相位为何与驱动快慢无关。
